\section{Step 60 Feed reconstruction from reconstructed coefficients}\label{h263-quant-s10}
The next connection is from the second output array to the inverse-transform
entry. For INTER, add each signed inverse residual to its predicted sample in a
wide signed value, then clip the sum to 0 through 255. For INTRA, clip the inverse
output directly to the sample range. Keep these final pixel operations outside
the quantizer: its reconstruction values are still transform coefficients.

Retain both arrays until coefficient coding and local reconstruction have consumed
them. Coding will read the selected levels; the local picture path will read
the reconstructed coefficients. Later prediction must use that locally reconstructed
picture, including the changes caused by quantization and final clipping.

\paragraph{Observe} A coefficient of 16 at $Q=8$ becomes reconstructed
coefficient 23. Feeding 16 to the inverse would bypass the very loss the decoder sees.
\Needspace{8\baselineskip}
\paragraph{Check} Which coefficient block must the local inverse consume?
\begin{itemize}
  \item[A.] The inverse-quantized block derived from the transmitted levels.
  \item[B.] The original forward block regardless of selected levels.
  \item[C.] The preceding raw picture's bytes reinterpreted as halfwords.
\end{itemize}
